% ---------------------------------------------------------------------
% Draft of Sec. VIII (Results) for the IEEE Access manuscript.
%
% Every number below is produced by a script in analysis/ and is marked
% with it in a comment, so each can be regenerated:
%   VIII-A  analysis/error_budget.py           (gate-passing counts)
%   VIII-B  analysis/onset_benchmark.py        (CV, detector comparison)
%   VIII-C  analysis/offset_accuracy_stats.py  (accuracy, per configuration)
%           analysis/offset_ge_variants.py     (ground-effect variants)
%   VIII-D  analysis/offset_accuracy_stats.py  (variance decomposition)
%   VIII-E  analysis/error_budget.py           (propagated budget)
%           analysis/vi_e_geometric_box.py     (bound at the flown excursions)
%           analysis/ramp_linearity.py         (ramp tracking residual)
%           analysis/baseline_bias.py          (baseline gain and sizes)
%           analysis/baseline_stat_ablation.py (C fitted vs pinned to zero)
%   VIII-F  analysis/moff_flight_check.py      (flown vs final pipeline)
%
% PLACEHOLDERS the author must fill -- marked \TODO below:
%   - the excitation protocol counts (runs per configuration per direction)
%   - the free-flight tracking metrics with and without compensation
%     (this file has the commanded offsets, not the flight performance)
%
% Cross-references assumed to exist in the host document; rename before
% compiling into the manuscript:
%   \eqref{eq:cosh}      the two-segment rate model
%   \eqref{eq:weight}    weight from the moments
%   \eqref{eq:offset}    offset per direction
%   \eqref{eq:pivotfree} pivot-free average
%   Sec.~\ref{sec:vie}   the deviation bound
% ---------------------------------------------------------------------

\newcommand{\TODO}[1]{\textcolor{red}{[\textbf{TODO:} #1]}}

\section{Results}
\label{sec:results}

\subsection{Excitation campaign and gating}

Five centre-of-mass configurations were produced by relocating ballast
within the planar offset box, giving ten independent case/axis
configurations once roll and pitch are counted separately. Each was
excited in both tip directions at commanded ramp rates of $0.10$,
$0.30$, $0.65$ and $1.20$~N\,m/s. \TODO{state the number of runs per
configuration per direction and the total attempted}

Runs enter the identification only through the gates of Sec.~VI-B: a
commanded-versus-realised rate error below $3\%$, a linearity residual
below $30$~mN\,m RMSE about the best-fit line, and a minimum
post-onset sample count. $140$ runs passed and are the basis of
everything below. The gates are deliberately blunt --- they reject
ramps that were stepped or aborted rather than tuning the fit --- and
no run was removed on the basis of its identified value.

\subsection{Onset repeatability and the choice of detector}

Within a configuration and direction, the identified onset moment
repeats to a coefficient of variation of $2.9\%$ on average and
$5.8\%$ at worst. That figure is the primary measure of the method's
precision, and it is what the deviation bound of Sec.~\ref{sec:vie}
predicts should be achievable.

Because a two-segment fit could in principle be finding structure that
a simpler detector would find too, six alternatives were run over the
identical gated set and carried through the identical inversion
(Table~\ref{tab:bench}). Two are variants of the proposed estimator
--- one using CAD rather than identified constants, one widening the
search box --- and four are model-agnostic: a per-run nonlinear least
squares, two change-point detectors, and a CUSUM.

\begin{table}[t]
\centering
\caption{Onset detectors over the same $140$ gated runs, each carried
through the same inversion. ``CV'' is the within-configuration
coefficient of variation of the onset moment; ``offset RMS'' is the
delivered quantity against the load cell. $\Delta|$err$|$ is the mean
paired change in $|$error$|$ against the proposed method, with an
exact sign test over the ten configurations.}
\label{tab:bench}
\small
\setlength{\tabcolsep}{4pt}
\begin{tabular}{@{}l rr rr rr@{}}
\hline
 & CV & CV & offset & mean & $\Delta|$err$|$ & $p$ \\
method & mean & worst & RMS & & & \\
 & [\%] & [\%] & [mm] & [mm] & [mm] & \\
\hline
COSH (proposed)      & $2.9$ & $5.8$  & $1.65$ & $-0.93$ & ---     & --- \\
COSH, CAD constants  & $4.4$ & $12.3$ & $1.64$ & $-0.86$ & $+0.00$ & $0.75$\\
COSH, wide box       & $2.9$ & $5.7$  & $1.64$ & $-0.95$ & $-0.02$ & $0.75$\\
per-run NLS          & $6.7$ & $19.8$ & $1.72$ & $-1.07$ & $+0.02$ & $0.34$\\
CPD, Gaussian cost   & $4.2$ & $8.6$  & $1.67$ & $-0.91$ & $-0.03$ & $0.75$\\
CPD, RBF kernel      & $9.3$ & $27.4$ & $1.65$ & $-0.46$ & $-0.09$ & $1.00$\\
CUSUM, mean change   & $7.7$ & $24.0$ & $1.82$ & $-0.58$ & $+0.29$ & $0.75$\\
\hline
\end{tabular}
\end{table}

Two things follow, and they point in opposite directions. The proposed
estimator is the most repeatable by a wide margin --- $2.9\%$ against
$4.2$ to $9.3\%$ --- which matters for a pre-flight check that must
return the same answer twice. But every method delivers an offset RMS
between $1.64$ and $1.82$~mm, no paired comparison approaches
significance ($p\ge0.34$), and the delivered offsets agree to
$1.09$~mm in the median across all seven. The onset is therefore
\emph{not} the limiting step: the accuracy of the deliverable is set
elsewhere, which Sec.~VIII-D takes up.

\subsection{Identified offset against the load cell}

The deliverable is the pivot-free average \eqref{eq:pivotfree},
inverted to a centre-of-mass offset and compared against an
independent load-cell measurement of the same quantity
(Table~\ref{tab:offset}).

\begin{table}[t]
\centering
\caption{Identified centre-of-mass offset against the load-cell truth,
over the ten configurations. ``repeat sd'' is the standard deviation
over run pairs within the configuration, and ``err/sd'' expresses the
error in units of that scatter.}
\label{tab:offset}
\small
\setlength{\tabcolsep}{4pt}
\begin{tabular}{@{}ll rrr rr@{}}
\hline
case & comp. & ident. & truth & error & repeat & err \\
 & & [mm] & [mm] & [mm] & sd [mm] & /sd \\
\hline
case\_01 & $y_{\rm off}$ & $-1.60$  & $-2.90$  & $+1.30$ & $1.09$ & $1.2$\\
case\_01 & $x_{\rm off}$ & $-11.18$ & $-11.45$ & $+0.27$ & $0.47$ & $0.6$\\
case\_02 & $y_{\rm off}$ & $-15.15$ & $-14.29$ & $-0.86$ & $0.40$ & $2.1$\\
case\_02 & $x_{\rm off}$ & $-10.57$ & $-9.90$  & $-0.67$ & $0.44$ & $1.5$\\
case\_03 & $y_{\rm off}$ & $-8.60$  & $-5.26$  & $-3.34$ & $0.51$ & $6.5$\\
case\_03 & $x_{\rm off}$ & $0.72$   & $3.14$   & $-2.42$ & $0.44$ & $5.4$\\
case\_04 & $y_{\rm off}$ & $5.24$   & $6.67$   & $-1.43$ & $0.74$ & $1.9$\\
case\_04 & $x_{\rm off}$ & $0.22$   & $2.40$   & $-2.18$ & $0.69$ & $3.2$\\
case\_05 & $y_{\rm off}$ & $10.79$  & $10.91$  & $-0.12$ & $0.82$ & $0.2$\\
case\_05 & $x_{\rm off}$ & $-10.57$ & $-10.89$ & $+0.32$ & $0.77$ & $0.4$\\
\hline
\end{tabular}
\end{table}

Over the ten configurations the RMS error is $1.64$~mm, with a
bootstrap $95\%$ confidence interval of $0.92$ to $2.24$~mm ($20$k
resamples), and the worst configuration is $3.34$~mm. Against offsets
that span $\pm15$~mm this is a relative accuracy of roughly one part
in ten of the range identified.

The mean error is $-0.91$~mm with a confidence interval of $-1.94$ to
$+0.11$~mm; seven of ten errors are negative, which an exact sign test
does not establish ($p = 0.34$). We therefore report the negative mean
as a point estimate and not as a demonstrated bias --- ten
configurations do not settle it, and the interval includes zero.

Two internal consistency checks accompany the comparison, neither of
which uses the load cell. Eq.~\eqref{eq:offset} estimates the same
offset twice, once from each tip direction with that direction's own
contact arm, and \eqref{eq:weight} returns the vehicle weight from the
onset moments. Without a ground-effect term the two directional
estimates disagree by $12.6$~mm in the median and \eqref{eq:weight}
returns a weight $5.2\%$ below $mg$ in all ten configurations;
restoring the interference channels brings the spread to $3.6$~mm and
the weight to $1.1\%$ below. The identified offset itself is reported
without a ground-effect model, to match the quantity that is flown;
Sec.~VI-B bounds what that omission costs at about half a millimetre
either way.

\subsection{What limits the accuracy}

The errors of Table~\ref{tab:offset} are not run scatter. Within a
configuration the run-pair standard deviation is $0.60$~mm in the
median and $1.09$~mm at worst, whereas the standard deviation of the
error \emph{between} configurations is $1.43$~mm --- a ratio of
$2.4$. Read against each configuration's own scatter, three of the ten
errors exceed three standard deviations, the worst reaching $6.5$.
The error is therefore systematic per configuration, and repeating a
configuration cannot reduce it. This is the central statistical
statement of the paper: the reported accuracy is limited by the
ground-contact geometry, not by the number of excitation runs.

\subsection{The realised error budget}
\label{sec:realised}

Sec.~\ref{sec:vie} bounds the estimator before any run, from the
operating box alone: $0.372$--$0.400$~mm for roll and
$0.304$--$0.331$~mm for pitch over the ramp-rate range, at the
$10^\circ$ tilt cap. Three refinements are available only after the
campaign, and they are collected here rather than there so that the
guarantee stays free of the data it is meant to bound.

\emph{The runs sit well inside the box.} Evaluating the same bound at
the excursion each window actually reached --- $5.95^\circ$ in the
median, $9.33^\circ$ at worst, with $x = C_2\tau_{\rm end}$ running
$1.79$ to $5.20$ --- gives $|\Delta\lambda_{\rm off}| \le 0.095$~mm in
the median and $0.249$~mm at worst. The excursions the protocol
actually produced are what buys the factor of four over the box.

\emph{The propagated error is smaller again.} The bound forwards the
unprojected channel maxima; the per-run propagation convolves the
\emph{measured} out-of-span forcing with the exact kernel after the
estimator has absorbed its own span. That gives a residual of
$2.75$~mN\,m in the median ($12.28$ at worst) on the gravity channel
and $0.53$ ($3.20$) on the ground-effect channel, a rate deviation of
$0.24\%$ of nominal in the median, and a threshold error of
$0.04$~mN\,m in the median and $0.15$ at worst --- $0.001$ to
$0.004$~mm of offset, some fifty times inside the bound. The
ground-effect channel supplies $5.1\%$ of $\bar\rho$ in the median and
$25.6\%$ at worst, but $0.00$--$0.01$~mN\,m of the threshold error once
propagated: it is not the channel that sets the budget.

\emph{The applied ramp is not exactly linear, and that is measured.}
The model presumes $M(\tau) = M_0 + \dot M\tau$, and $\dot M$ is the
least-squares slope of the \emph{measured} moment over the window, so
the constant and linear parts of any tracking error are absorbed by the
normal equations and only the non-affine part survives. Across the
$140$ runs that residual is $2.01$~mN\,m RMS in the median and
$6.49$~mN\,m peak ($23.96$ at worst), i.e.\ $0.5$--$0.8\%$ of the
moment span; the worst run's RMS is $4.43$~mN\,m, $6.8$ times inside
the $30$~mN\,m linearity gate of Sec.~VI-B. Because it is orthogonal to $\{1,\tau\}$ by construction, it
can only pair with the non-affine part of the reduction weight, which
carries $\|w - P_1w\|_1 = 0.54$ of the weight's mass at the largest $x$
in the box; the resulting threshold contribution is $0.117$~mm in the
median and $0.431$~mm at worst. This channel lies outside the a priori
bound, and at its worst it is the largest of the three --- still an
order of magnitude below the $1.64$~mm achieved.

\emph{The pre-onset baseline is estimated, not assumed.} Eq.~(106)
shows the identified threshold responding to the baseline \emph{error}
$C - b$ with a gain of $0.042$ to $0.351$~N\,m per rad/s over the box.
The baselines the sweep forms --- medians over the excitation window up
to the candidate --- are $3.5$~mrad/s across the $140$ runs and
$47.8$~mrad/s at worst, which bounds the term at $0.054$~mm in the
median and $0.558$~mm at worst. That bound is loose, and the direct
test says by how much: re-running the whole pipeline with $C$ pinned to
zero, the largest perturbation the term admits and the one the onset
conditions would motivate, moves the direction-paired offset by
$0.027$~mm in the median and $0.163$~mm at worst. The pairing cancels
most of what the bound forwards. Neither choice measures better --- the
ramp-invariance coefficient of variation is $0.0287$ fitted against
$0.0292$ pinned, a paired $t$ of $0.87$ on $19$ degrees of freedom, and
$66$ of the $140$ runs return an identical onset index either way. $C$
is fitted, and fitted as a median, because a pinned zero has no defence
against genuine pre-excitation motion, not because it scores better.

Every figure above is at least four times inside that $1.64$~mm, which
locates the residual outside the estimator and is consistent with the
detector comparison of Table~\ref{tab:bench}.

\subsection{Free-flight compensation}

The identified offsets were applied as a feed-forward moment trim at
take-off. The values flown are compared in Table~\ref{tab:flown}
against what the final pipeline returns, since the pipeline was
refined after the flights were performed.

\begin{table}[t]
\centering
\caption{Moment offsets commanded in the free-flight tests against the
final pipeline. ``std'' is the scatter of the identification at the
time of flight.}
\label{tab:flown}
\small
\setlength{\tabcolsep}{4pt}
\begin{tabular}{@{}ll rrr r@{}}
\hline
case & comp. & flown & std & final & diff \\
 & & [N\,m] & [N\,m] & [N\,m] & /std \\
\hline
case\_01 & $M_{{\rm off},x}$ & $-0.064$ & $0.037$ & $-0.048$ & $0.43$\\
case\_01 & $M_{{\rm off},y}$ & $0.331$  & $0.016$ & $0.336$  & $0.33$\\
case\_02 & $M_{{\rm off},x}$ & $-0.485$ & $0.009$ & $-0.479$ & $0.71$\\
case\_02 & $M_{{\rm off},y}$ & $0.333$  & $0.010$ & $0.334$  & $0.10$\\
case\_03 & $M_{{\rm off},x}$ & $-0.282$ & $0.028$ & $-0.272$ & $0.37$\\
case\_03 & $M_{{\rm off},y}$ & $-0.006$ & $0.019$ & $-0.023$ & $0.89$\\
case\_04 & $M_{{\rm off},x}$ & $0.173$  & $0.031$ & $0.165$  & $0.24$\\
case\_04 & $M_{{\rm off},y}$ & $-0.012$ & $0.024$ & $-0.007$ & $0.21$\\
case\_05 & $M_{{\rm off},x}$ & $0.307$  & $0.026$ & $0.341$  & $1.30$\\
case\_05 & $M_{{\rm off},y}$ & $0.343$  & $0.031$ & $0.334$  & $0.30$\\
\hline
\end{tabular}
\end{table}

The two agree to $8$~mN\,m in the median and $34$~mN\,m at worst, and
nine of the ten components fall inside the identification scatter that
was present at the time of flight; the worst is $1.3$ standard
deviations. Re-flying would therefore command the same moments to
within the precision of the identification itself, and the free-flight
results are reported against the pipeline as described.

\TODO{Report the flight performance here: attitude tracking error and
take-off transient with and without compensation, three repetitions per
configuration per condition, with the paired statistic across the five
configurations. This file supplies the commanded offsets only.}

\subsection{Summary}

The method identifies a centre-of-mass offset to $1.64$~mm RMS
($95\%$ CI $0.92$--$2.24$) without a load cell, a test stand, or any
flight, from ground-contact excitation alone. It repeats to $2.9\%$ on
the onset moment, better than every model-agnostic baseline tried,
although the delivered offset is insensitive to that choice. Its own
modelling error is bounded a priori at $0.174$~mm and measured at
$0.001$--$0.004$~mm, so what limits the accuracy is the contact
geometry rather than the estimator or the number of repetitions ---
a conclusion the variance decomposition establishes directly, and one
that determines where further work should go.
